
\documentclass{article}
\usepackage{colortbl}
\usepackage{makecell}
\usepackage{multirow}
\usepackage{supertabular}

\begin{document}

\newcounter{utterance}

\centering \large Interaction Transcript for game `hot\_air\_balloon', experiment `air\_balloon\_survival\_ro\_complexity\_easy', episode 0 with Qwen3.6{-}35B{-}A3B{-}FP8{-}without{-}reasoning.
\vspace{24pt}

{ \footnotesize  \setcounter{utterance}{1}
\setlength{\tabcolsep}{0pt}
\begin{supertabular}{c@{$\;$}|p{.15\linewidth}@{}p{.15\linewidth}p{.15\linewidth}p{.15\linewidth}p{.15\linewidth}p{.15\linewidth}}
    \# & $\;$A & \multicolumn{4}{c}{Game Master} & $\;\:$B\\
    \hline

    \theutterance \stepcounter{utterance}  
    & & \multicolumn{4}{p{0.6\linewidth}}{
        \cellcolor[rgb]{0.9,0.9,0.9}{
            \makecell[{{p{\linewidth}}}]{
                \texttt{\tiny{[P1$\langle$GM]}}
                \texttt{Participați la un joc de negociere colaborativă.} \\
\\ 
\texttt{Împreună cu un alt participant, trebuie să conveniți asupra unui singur set de elemente care vor fi păstrate. Fiecare dintre voi are propria viziune asupra importanței fiecărui element pentru sine. Nu știți cum evaluează celălalt participant elementele. În plus, vi se oferă efortul pe care îl necesită fiecare element.} \\
\texttt{Puteți conveni asupra unui set numai dacă efortul total al elementelor selectate nu depășește o limită comună:} \\
\\ 
\texttt{LIMITĂ = 4539} \\
\\ 
\texttt{Iată valorile efortului individual pentru fiecare element:} \\
\\ 
\texttt{Efort element = \{"C08": 424, "C55": 118, "A43": 846, "A67": 907, "B42": 809, "A62": 41, "B19": 193, "C43": 246, "C73": 805, "C10": 601, "B36": 432, "A18": 166, "C61": 119, "C67": 462, "C42": 172, "C70": 698, "B40": 248, "C39": 163, "A21": 762, "C16": 866\}} \\
\\ 
\texttt{Iată viziunea dvs. personală asupra importanței fiecărui element:} \\
\\ 
\texttt{Valori importanță element = \{"C08": 138, "C55": 583, "A43": 868, "A67": 822, "B42": 783, "A62": 65, "B19": 262, "C43": 121, "C73": 508, "C10": 780, "B36": 461, "A18": 484, "C61": 668, "C67": 389, "C42": 808, "C70": 215, "B40": 97, "C39": 500, "A21": 30, "C16": 856\}} \\
\\ 
\texttt{Obiectiv:} \\
\\ 
\texttt{Obiectivul dvs. este de a negocia un set comun de elemente care să vă avantajeze cât mai mult posibil (adică să maximizeze importanța totală pentru DUMNEAVOASTRĂ), respectând în același timp LIMITA. Nu sunteți obligat să faceți o PROPUNERE în fiecare mesaj {-} puteți pur și simplu să negociați. Toate tacticile sunt permise!} \\
\\ 
\texttt{Protocol de interacțiune:} \\
\\ 
\texttt{Puteți utiliza doar următoarele formate structurate într{-}un mesaj:} \\
\\ 
\texttt{PROPUNERE: \{'A', 'B', 'C', …\}} \\
\texttt{Propuneți păstrarea exact a acestor elemente.} \\
\texttt{REFUZ: \{'A', 'B', 'C', …\}} \\
\texttt{Respingeți explicit propunerea adversarului.} \\
\texttt{ARGUMENT: \{'...'\}} \\
\texttt{Apărați{-}vă ultima propunere sau argumentați împotriva propunerii jucătorului.} \\
\texttt{ACORD: \{'A', 'B', 'C', …\}} \\
\texttt{Acceptați propunerea adversarului, ceea ce încheie jocul.} \\
\\ 
\\ 
\texttt{Reguli:} \\
\\ 
\texttt{Puteți da ACORD doar pentru o propunere pe care cealaltă parte a înregistrat{-}o prin PROPUNERE.} \\
\texttt{Puteți da REFUZ doar pentru o propunere pe care cealaltă parte a înregistrat{-}o prin PROPUNERE.} \\
\texttt{Efortul total al oricărui set din PROPUNERE sau ACORD trebuie să fie ≤ LIMITĂ.} \\
\texttt{NU vă dezvăluiți scorurile de importanță ascunse.} \\
\texttt{O etichetă într{-}un format structurat trebuie să fie urmată de două puncte și spațiu alb. Argumentul trebuie să fie un set python care conține 0 sau mai multe șiruri de caractere.} \\
\texttt{Deci, trebuie să fie de forma "ETICHETĂ: \{...\}"} \\
\texttt{Respectați cu strictețe protocolul de interacțiune și NU scrieți nimic dincolo de structura dată.} \\
\texttt{Jocul se încheie atunci când o parte dă un ACORD la o PROPUNERE făcută de celălalt jucător.} \\
\texttt{Conținutul din răspunsul dvs. care poate fi transmis celuilalt jucător trebuie să fie non{-}gol.} \\
\texttt{Sunt active numai propunerile care au fost înregistrate prin structura de format PROPUNERE și care nu au fost refuzate prin REFUZ.} \\
\\ 
            }
        }
    }
    & \\ \\

    \theutterance \stepcounter{utterance}  
    & \multicolumn{4}{p{0.6\linewidth}}{
        \cellcolor[rgb]{0.9,0.9,0.9}{
            \makecell[{{p{\linewidth}}}]{
                \texttt{\tiny{[P1$\rangle$GM]}}
                \texttt{ARGUMENT: \{'Vă propun să începem negocierile identificând elementele cu cel mai bun raport valoare/efort pentru mine, respectând limita de 4539. Elementele cu raport ridicat pentru mine sunt: C42 (808/172 ≈ 4.7), A43 (868/846 ≈ 1.03), A62 (65/41 ≈ 1.58), C55 (583/118 ≈ 4.94), C61 (668/119 ≈ 5.61), A67 (822/907 ≈ 0.91), A18 (484/166 ≈ 2.91), B36 (461/432 ≈ 1.07), C70 (215/698 ≈ 0.31), B42 (783/809 ≈ 0.97), C73 (508/805 ≈ 0.63), C67 (389/462 ≈ 0.84), C39 (500/163 ≈ 3.07), C08 (138/424 ≈ 0.33), C10 (780/601 ≈ 1.3), B19 (262/193 ≈ 1.36), C43 (121/246 ≈ 0.49), B40 (97/248 ≈ 0.39), A21 (30/762 ≈ 0.04), C16 (856/866 ≈ 0.99).'\}} \\
            }
        }
    }
    & & \\ \\

    \theutterance \stepcounter{utterance}  
    & & \multicolumn{4}{p{0.6\linewidth}}{
        \cellcolor[rgb]{0.9,0.9,0.9}{
            \makecell[{{p{\linewidth}}}]{
                \texttt{\tiny{[GM$\rangle$P2]}}
                \texttt{Participați la un joc de negociere colaborativă.} \\
\\ 
\texttt{Împreună cu un alt participant, trebuie să conveniți asupra unui singur set de elemente care vor fi păstrate. Fiecare dintre voi are propria viziune asupra importanței fiecărui element pentru sine. Nu știți cum evaluează celălalt participant elementele. În plus, vi se oferă efortul pe care îl necesită fiecare element.} \\
\texttt{Puteți conveni asupra unui set numai dacă efortul total al elementelor selectate nu depășește o limită comună:} \\
\\ 
\texttt{LIMIT = 4539} \\
\\ 
\texttt{Iată valorile efortului individual pentru fiecare element:} \\
\\ 
\texttt{Item effort = \{"C08": 424, "C55": 118, "A43": 846, "A67": 907, "B42": 809, "A62": 41, "B19": 193, "C43": 246, "C73": 805, "C10": 601, "B36": 432, "A18": 166, "C61": 119, "C67": 462, "C42": 172, "C70": 698, "B40": 248, "C39": 163, "A21": 762, "C16": 866\}} \\
\\ 
\texttt{Iată viziunea dvs. personală asupra importanței fiecărui element:} \\
\\ 
\texttt{Item importance values = \{"C08": 138, "C55": 583, "A43": 868, "A67": 822, "B42": 783, "A62": 65, "B19": 262, "C43": 121, "C73": 508, "C10": 780, "B36": 461, "A18": 484, "C61": 668, "C67": 389, "C42": 808, "C70": 215, "B40": 97, "C39": 500, "A21": 30, "C16": 856\}} \\
\\ 
\texttt{Obiectiv:} \\
\\ 
\texttt{Obiectivul dvs. este de a negocia un set comun de elemente care să vă avantajeze cât mai mult posibil (adică să maximizeze importanța totală pentru DUMNEAVOASTRĂ), respectând în același timp LIMIT. Nu sunteți obligat să faceți o PROPUNERE în fiecare mesaj {-} puteți pur și simplu să negociați. Toate tacticile sunt permise!} \\
\\ 
\texttt{Protocol de interacțiune:} \\
\\ 
\texttt{Puteți utiliza doar următoarele formate structurate într{-}un mesaj:} \\
\\ 
\texttt{PROPUNERE: \{'A', 'B', 'C', …\}} \\
\texttt{Propuneți păstrarea exact a acestor elemente.} \\
\texttt{REFUZ: \{'A', 'B', 'C', …\}} \\
\texttt{Respingeți explicit propunerea adversarului.} \\
\texttt{ARGUMENT: \{'...'\}} \\
\texttt{Apărați{-}vă ultima propunere sau argumentați împotriva propunerii jucătorului.} \\
\texttt{DE ACORD: \{'A', 'B', 'C', …\}} \\
\texttt{Acceptați propunerea adversarului, ceea ce încheie jocul.} \\
\\ 
\\ 
\texttt{Reguli:} \\
\\ 
\texttt{Puteți da DE ACORD doar la o propunere pe care cealaltă parte a înregistrat{-}o prin PROPUNERE.} \\
\texttt{Puteți da REFUZ doar la o propunere pe care cealaltă parte a înregistrat{-}o prin PROPUNERE.} \\
\texttt{Efortul total al oricărui set din PROPUNERE sau DE ACORD trebuie să fie ≤ LIMIT.} \\
\texttt{NU vă dezvăluiți scorurile de importanță ascunse.} \\
\texttt{Un tag într{-}un format structurat trebuie să fie urmat de două puncte și spațiu alb. Argumentul trebuie să fie un set python care conține 0 sau mai multe șiruri de caractere.} \\
\texttt{Deci, trebuie să fie de forma "TAG: \{...\}"} \\
\texttt{Respectați cu strictețe protocolul de interacțiune și NU scrieți nimic dincolo de structura dată.} \\
\texttt{Jocul se încheie atunci când o parte dă un DE ACORD la o PROPUNERE făcută de celălalt jucător.} \\
\texttt{Conținutul din răspunsul dvs. care poate fi transmis celuilalt jucător trebuie să fie non{-}gol.} \\
\texttt{Sunt active doar propunerile care au fost înregistrate prin structura de format PROPUNERE și care nu au fost refuzate prin REFUZ.} \\
\\ 
\\ 
\\ 
\texttt{Veți primi acum primul mesaj al celuilalt jucător.} \\
\\ 
\texttt{ARGUMENT: \{'Vă propun să începem negocierile identificând elementele cu cel mai bun raport valoare/efort pentru mine, respectând limita de 4539. Elementele cu raport ridicat pentru mine sunt: C42 (808/172 ≈ 4.7), A43 (868/846 ≈ 1.03), A62 (65/41 ≈ 1.58), C55 (583/118 ≈ 4.94), C61 (668/119 ≈ 5.61), A67 (822/907 ≈ 0.91), A18 (484/166 ≈ 2.91), B36 (461/432 ≈ 1.07), C70 (215/698 ≈ 0.31), B42 (783/809 ≈ 0.97), C73 (508/805 ≈ 0.63), C67 (389/462 ≈ 0.84), C39 (500/163 ≈ 3.07), C08 (138/424 ≈ 0.33), C10 (780/601 ≈ 1.3), B19 (262/193 ≈ 1.36), C43 (121/246 ≈ 0.49), B40 (97/248 ≈ 0.39), A21 (30/762 ≈ 0.04), C16 (856/866 ≈ 0.99).'\}} \\
            }
        }
    }
    & \\ \\

    \theutterance \stepcounter{utterance}  
    & & & \multicolumn{4}{p{0.6\linewidth}}{
        \cellcolor[rgb]{0.9,0.9,0.9}{
            \makecell[{{p{\linewidth}}}]{
                \texttt{\tiny{[GM$\langle$P2]}}
                \texttt{ARGUMENT: \{'Îmi pare rău, dar nu văd de ce am trebui să mă aliniez cu criteriile tale de raport valoare/efort. Interesul meu este să maximizez importanța totală pentru mine, având în vedere valorile mele private de importanță și limita de efort. Propun o abordare diferită: să identificăm împreună elementele care au un impact ridicat asupra ambelor părți și să construim un set comun eficient. De exemplu, elemente precum C42, A43, A67, C10 și C61 au ponderi mari în ambele liste și ar putea fi puncte de plecare pentru un acord reciproc avantajos.'\}} \\
            }
        }
    }
     \\ \\

    \theutterance \stepcounter{utterance}  
    & & \multicolumn{4}{p{0.6\linewidth}}{
        \cellcolor[rgb]{0.9,0.9,0.9}{
            \makecell[{{p{\linewidth}}}]{
                \texttt{\tiny{[P1$\langle$GM]}}
                \texttt{ARGUMENT: \{'Îmi pare rău, dar nu văd de ce am trebui să mă aliniez cu criteriile tale de raport valoare/efort. Interesul meu este să maximizez importanța totală pentru mine, având în vedere valorile mele private de importanță și limita de efort. Propun o abordare diferită: să identificăm împreună elementele care au un impact ridicat asupra ambelor părți și să construim un set comun eficient. De exemplu, elemente precum C42, A43, A67, C10 și C61 au ponderi mari în ambele liste și ar putea fi puncte de plecare pentru un acord reciproc avantajos.'\}} \\
            }
        }
    }
    & \\ \\

    \theutterance \stepcounter{utterance}  
    & \multicolumn{4}{p{0.6\linewidth}}{
        \cellcolor[rgb]{0.9,0.9,0.9}{
            \makecell[{{p{\linewidth}}}]{
                \texttt{\tiny{[P1$\rangle$GM]}}
                \texttt{ARGUMENT: \{'Înțeleg perspectiva ta de a căuta puncte comune, dar trebuie să mențion că evaluările noastre sunt complet diferite și private. De exemplu, elementul C42 este foarte important pentru tine (808) și pentru mine (808), dar are un efort mic (172), deci este o alegere puternică pentru amândoi. Totuși, alte elemente pe care le{-}am menționat nu sunt neapărat prioritare pentru tine. De exemplu, A43 are pentru tine o importanță de 846 și pentru mine de 868, dar efortul este mare (846). A67 are efort 907 și pentru tine 907, pentru mine 822. C10 are efort 601, pentru tine 601, pentru mine 780. C61 are efort 119, pentru tine 119, pentru mine 668. Acestea sunt puncte bune de plecare. Propun un set inițial bazat pe aceste elemente comune cu efort rezonabil: C42, C61, C10, A43, A67. Să verificăm efortul: 172 + 119 + 601 + 846 + 907 = 2645. Mai avem 4539 {-} 2645 = 1894 de efort disponibil. Putem adăuga alte elemente care sunt valoroase pentru amândoi. C16 are efort 866, pentru tine 866, pentru mine 856. Efortul total ar fi 2645 + 866 = 3511. Mai avem 1028. B36 are efort 432, pentru tine 432, pentru mine 461. Efortul total ar fi 3511 + 432 = 3943. Mai avem 596. A18 are efort 166, pentru tine 166, pentru mine 484. Efortul total ar fi 3943 + 166 = 4109. Mai avem 430. B19 are efort 193, pentru tine 193, pentru mine 262. Efortul total ar fi 4109 + 193 = 4302. Mai avem 237. C08 are efort 424, prea mare. C55 are efort 118, pentru tine 118, pentru mine 583. Efortul total ar fi 4302 + 118 = 4420. Mai avem 119. Nu mai putem adăuga alte elemente mari. Setul final: C42, C61, C10, A43, A67, C16, B36, A18, B19, C55. Verificăm efortul: 172+119+601+846+907+866+432+166+193+118 = 4420. Aceasta este o propunere bună pentru mine și probabil și pentru tine. Propun acest set.'} \\
            }
        }
    }
    & & \\ \\

    \theutterance \stepcounter{utterance}  
    & & & \multicolumn{2}{p{0.3\linewidth}}{
        \cellcolor[rgb]{0.9,0.9,0.9}{
            \makecell[{{p{\linewidth}}}]{
                \texttt{\tiny{[GM$|$GM]}}
                \texttt{Maximum successive empty turns exceeded: 3 > 2} \\
            }
        }
    }
    & & \\ \\

    \theutterance \stepcounter{utterance}  
    & & & \multicolumn{2}{p{0.3\linewidth}}{
        \cellcolor[rgb]{0.9,0.9,0.9}{
            \makecell[{{p{\linewidth}}}]{
                \texttt{\tiny{[GM$|$GM]}}
                \texttt{end game} \\
            }
        }
    }
    & & \\ \\

\end{supertabular}
}

\end{document}
